\documentclass[10pt,a4paper]{amsart}
\usepackage{amsmath,amssymb,amsthm,graphicx}
% Retained minimal preamble for the Stacks Project chapter sites-cohomology.tex (Cohomology on Sites).
% Verbatim excerpt of the chapter preamble: Xy-pic setup (native preamble lines 14-18) and the native
% theorem declarations with the mathematical macros (native preamble lines 154-217). Copied mechanically;
% nothing added, removed or redefined. The wrapper class and the amsmath/amssymb/amsthm/graphicx base
% packages are supplied by the assembler, so the class/externaldocument part of the native preamble is omitted.
% For commutative diagrams we use Xy-pic
\usepackage[all]{xy}

% We use 2cell for 2-commutative diagrams.
\xyoption{2cell}
% Theorem environments.
%
\theoremstyle{plain}
\newtheorem{theorem}[subsection]{Theorem}
\newtheorem{proposition}[subsection]{Proposition}
\newtheorem{lemma}[subsection]{Lemma}

\theoremstyle{definition}
\newtheorem{definition}[subsection]{Definition}
\newtheorem{example}[subsection]{Example}
\newtheorem{exercise}[subsection]{Exercise}
\newtheorem{situation}[subsection]{Situation}

\theoremstyle{remark}
\newtheorem{remark}[subsection]{Remark}
\newtheorem{remarks}[subsection]{Remarks}

\numberwithin{equation}{subsection}

% Macros
%
\def\lim{\mathop{\mathrm{lim}}\nolimits}
\def\colim{\mathop{\mathrm{colim}}\nolimits}
\def\Spec{\mathop{\mathrm{Spec}}}
\def\Hom{\mathop{\mathrm{Hom}}\nolimits}
\def\Ext{\mathop{\mathrm{Ext}}\nolimits}
\def\SheafHom{\mathop{\mathcal{H}\!\mathit{om}}\nolimits}
\def\SheafExt{\mathop{\mathcal{E}\!\mathit{xt}}\nolimits}
\def\Sch{\mathit{Sch}}
\def\Mor{\mathop{\mathrm{Mor}}\nolimits}
\def\Ob{\mathop{\mathrm{Ob}}\nolimits}
\def\Sh{\mathop{\mathit{Sh}}\nolimits}
\def\NL{\mathop{N\!L}\nolimits}
\def\CH{\mathop{\mathrm{CH}}\nolimits}
\def\proetale{{pro\text{-}\acute{e}tale}}
\def\etale{{\acute{e}tale}}
\def\QCoh{\mathit{QCoh}}
\def\Ker{\mathop{\mathrm{Ker}}}
\def\Im{\mathop{\mathrm{Im}}}
\def\Coker{\mathop{\mathrm{Coker}}}
\def\Coim{\mathop{\mathrm{Coim}}}

% Boxtimes
%
\DeclareMathSymbol{\boxtimes}{\mathbin}{AMSa}{"02}

%
% Macros for moduli stacks/spaces
%
\def\QCohstack{\mathcal{QC}\!\mathit{oh}}
\def\Cohstack{\mathcal{C}\!\mathit{oh}}
\def\Spacesstack{\mathcal{S}\!\mathit{paces}}
\def\Quotfunctor{\mathrm{Quot}}
\def\Hilbfunctor{\mathrm{Hilb}}
\def\Curvesstack{\mathcal{C}\!\mathit{urves}}
\def\Polarizedstack{\mathcal{P}\!\mathit{olarized}}
\def\Complexesstack{\mathcal{C}\!\mathit{omplexes}}
% \Pic is the operator that assigns to X its picard group, usage \Pic(X)
% \Picardstack_{X/B} denotes the Picard stack of X over B
% \Picardfunctor_{X/B} denotes the Picard functor of X over B
\def\Pic{\mathop{\mathrm{Pic}}\nolimits}
\def\Picardstack{\mathcal{P}\!\mathit{ic}}
\def\Picardfunctor{\mathrm{Pic}}
\def\Deformationcategory{\mathcal{D}\!\mathit{ef}}
\usepackage{cleveref}
\crefname{originalreference}{original reference}{original references}
\Crefname{originalreference}{Original reference}{Original references}
\makeatletter
\newlabel{definition-K-flat}{{06XZ}{source}{}{sourcepointer-1}{}}
\newlabel{definition-K-flat@cref}{{[originalreference][0][]06XZ}{[1][0][]source}{}{}{}}
\makeatother
\crefname{originalequation}{original equation}{original equations}
\Crefname{originalequation}{Original equation}{Original equations}
\makeatletter
\newlabel{equation-global-cohomology-hom-complex}{{0A8Z}{source}{}{sourcepointer-2}{}}
\newlabel{equation-global-cohomology-hom-complex@cref}{{[originalequation][0][]0A8Z}{[1][0][]source}{}{}{}}
\makeatother
\crefname{originalreference}{original reference}{original references}
\Crefname{originalreference}{Original reference}{Original references}
\makeatletter
\newlabel{lemma-compose}{{0A5Y}{source}{}{sourcepointer-3}{}}
\newlabel{lemma-compose@cref}{{[originalreference][0][]0A5Y}{[1][0][]source}{}{}{}}
\makeatother
\title{The internal hom from a K-flat complex into a K-injective complex is again K-injective}
\author{The Stacks Project Authors}
\date{stacks-project:sites-cohomology.tex:local-snapshot-20261011; GNU FDL 1.2 or later}
\begin{document}
\maketitle
\noindent\textit{Editorial note.} All delivered mathematical text is the author's own, copied from the pinned Stacks Project chapter file sites-cohomology.tex; no statement or proof step was rewritten, reordered or removed.\par\medskip
\noindent\textit{Editorial note.} The counted proof is the author's complete argument: the proof is written out in full: the displayed four-term chain of equalities with the reason for each, and the final vanishing from K-flatness of the tensor factor and K-injectivity of the target.\par\medskip
\noindent\textit{Editorial note.} The label lemma-RHom-from-K-flat-into-K-injective names the counted statement, and the proof environment immediately following it in the source is delivered verbatim as its proof.\par\medskip
\noindent\textit{Editorial note.} The counted statement is carried by the author's own lemma environment in the chapter source.\par\medskip
\noindent\textit{Editorial note.} The statement and proof are copied from the pinned chapter file; the three cited inputs are kept as external pointers with their canonical Stacks tags.\par\medskip
\noindent\textit{Editorial note.} The preamble is the chapter's own retained minimal preamble (sources/stacks-sites-cohomology/retained-preamble.tex), a verbatim excerpt of the native preamble: the Xy-pic setup and the native theorem declarations with the mathematical macros.\par\medskip
\noindent\textit{Allowed standard prerequisites (editorial).}\begin{itemize}
\item K-flat and K-injective complexes of modules on a ringed site. Same chapter, the definitions of K-flat and K-injective complexes used as the two hypotheses of the statement. Supplies the properties invoked at the end of the proof: the total tensor product with a K-flat complex of an acyclic complex is acyclic, and a K-injective complex has no homotopy classes of maps from acyclic complexes.
\item The global cohomology description of Hom complexes. Same chapter, the equation identifying Hom in the homotopy category with H\textasciicircum{}0 of the global sections of the internal hom complex. Supplies the first and third equalities of the displayed chain.
\item Composition of internal homs and the total tensor product. Same chapter, the lemma identifying the internal hom into an internal hom with the internal hom from the total tensor product. Supplies the second equality of the displayed chain.
\end{itemize}\medskip
\ifdefined\hypertarget\hypertarget{sourcepointer-1}{}\fi
\noindent\textit{Source reference.} Original reference 06XZ: sources/stacks-sites-cohomology/pinned-original-source.tex, source lines 2999--2999. Prerequisite pointer invoked at proof line(s) 8824; it names a standard fact of the same chapter that the counted argument uses, and every step of the counted proof built on it is retained verbatim in the delivered body.\par\medskip
\ifdefined\hypertarget\hypertarget{sourcepointer-2}{}\fi
\noindent\textit{Source reference.} Original equation 0A8Z: sources/stacks-sites-cohomology/pinned-original-source.tex, source lines 8608--8608. Prerequisite pointer invoked at proof line(s) 8818; it names a standard fact of the same chapter that the counted argument uses, and every step of the counted proof built on it is retained verbatim in the delivered body.\par\medskip
\ifdefined\hypertarget\hypertarget{sourcepointer-3}{}\fi
\noindent\textit{Source reference.} Original reference 0A5Y: sources/stacks-sites-cohomology/pinned-original-source.tex, source lines 8616--8616. Prerequisite pointer invoked at proof line(s) 8819; it names a standard fact of the same chapter that the counted argument uses, and every step of the counted proof built on it is retained verbatim in the delivered body.\par\medskip
\begin{lemma}
\label{lemma-RHom-from-K-flat-into-K-injective}
Let $(\mathcal{C}, \mathcal{O})$ be a ringed site. Let $\mathcal{I}^\bullet$
be a K-injective complex of $\mathcal{O}$-modules. Let
$\mathcal{L}^\bullet$ be a K-flat complex of $\mathcal{O}$-modules.
Then $\SheafHom^\bullet(\mathcal{L}^\bullet, \mathcal{I}^\bullet)$
is a K-injective complex of $\mathcal{O}$-modules.
\end{lemma}

\begin{proof}
Namely, if $\mathcal{K}^\bullet$ is an acyclic complex of
$\mathcal{O}$-modules, then
\begin{align*}
\Hom_{K(\mathcal{O})}(\mathcal{K}^\bullet,
\SheafHom^\bullet(\mathcal{L}^\bullet, \mathcal{I}^\bullet))
& =
H^0(\Gamma(\mathcal{C},
\SheafHom^\bullet(\mathcal{K}^\bullet,
\SheafHom^\bullet(\mathcal{L}^\bullet, \mathcal{I}^\bullet)))) \\
& =
H^0(\Gamma(\mathcal{C}, \SheafHom^\bullet(\text{Tot}(
\mathcal{K}^\bullet \otimes_\mathcal{O} \mathcal{L}^\bullet),
\mathcal{I}^\bullet))) \\
& =
\Hom_{K(\mathcal{O})}(
\text{Tot}(\mathcal{K}^\bullet \otimes_\mathcal{O} \mathcal{L}^\bullet),
\mathcal{I}^\bullet) \\
& =
0
\end{align*}
The first equality by (\ref{equation-global-cohomology-hom-complex}).
The second equality by Lemma \ref{lemma-compose}.
The third equality by (\ref{equation-global-cohomology-hom-complex}).
The final equality because
$\text{Tot}(\mathcal{K}^\bullet \otimes_\mathcal{O} \mathcal{L}^\bullet)$
is acyclic because $\mathcal{L}^\bullet$ is K-flat
(Definition \ref{definition-K-flat}) and because $\mathcal{I}^\bullet$
is K-injective.
\end{proof}

\end{document}
